% ============================================================================
\section{Instruction Set Principles}
\label{sec:isa}

This section covers everything on instruction sets in the course: ISA styles,
instruction set design (operations, addressing modes, operand types, formats,
registers), CISC vs.\ RISC, and the MIPS case study (formats, encoding,
addressing, constants, procedures, atomic operations and the assembler).

% ----------------------------------------------------------------------------
\subsection{What to know}

\subsubsection{The instruction set}
\begin{itemize}[leftmargin=*]
  \item The \textbf{instruction set} is the set of all instructions a
  processor can perform: the \textbf{interface between software and
  hardware}.
  \item Execution cycle: \textbf{instruction fetch} $\to$ \textbf{decode} $\to$
  \textbf{operand fetch} $\to$ \textbf{execute} (update status) $\to$
  \textbf{write back} $\to$ next fetch.
\end{itemize}

\subsubsection{ISA styles: where the operands live}
\begin{center}
\small
\begin{tabular}{@{}p{2.4cm}p{3.3cm}p{3.6cm}p{5.1cm}@{}}
\toprule
Style & $C=A+B$ & Operands & Pros / cons \\
\midrule
Stack & \texttt{PUSH A; PUSH B; ADD; POP C} & implicit, top of stack &
(+) code density, simple hardware; (--) low parallelism, stack is a bottleneck \\
Accumulator & \texttt{LOAD A; ADD B; STORE C} & one implicit register (AC) + memory &
(+) easy compiler, few instructions; (--) very high memory traffic, variable CPI \\
Memory-memory & \texttt{ADD C, A, B} & all in memory &
(+) simple to design and understand; (--) bottleneck, heavy memory access \\
Register-memory & \texttt{LOAD R1,A; ADD R3,R1,B; STORE R3,C} & register or memory &
(+) fewer instructions than load-store; (--) operand access time varies \\
Load-store (reg-reg) & \texttt{LOAD R1,A; LOAD R2,B; ADD R3,R1,R2; STORE R3,C} &
registers only; only loads/stores touch memory &
(+) simple fixed encoding, easy to pipeline, registers reused; (--) more
instructions \\
\bottomrule
\end{tabular}
\end{center}
Historical machines: Motorola 6800 (2 GPRs, accumulator, 1974), DEC VAX (16,
register-memory/memory-memory, 1977), Intel 8086 (1, extended accumulator,
1978), Motorola 68000 (16, register-memory, 1980), Intel 80386 (32,
register-memory, 1985), PowerPC and DEC Alpha (32, load-store, 1992). Other
styles: high-level-language architectures (Burroughs B5000, failed for lack of
compilers) and RISC.

\subsubsection{Instruction set design}
Five aspects: \textbf{operation repertoire}, \textbf{addressing modes},
\textbf{data types}, \textbf{instruction format}, \textbf{registers}. A simple
instruction has an \emph{opcode}, \emph{source operand reference(s)} and a
\emph{result operand reference}.
\begin{itemize}[leftmargin=*]
  \item \textbf{Instruction types} (4): data processing (arithmetic, logic),
  data storage (load/store to memory), data movement (register/register, I/O,
  \texttt{MOV}), control (test and branch: \texttt{CMP}, \texttt{JMP}).
  \item \textbf{Operator types} (8): arithmetic \& logical, data transfer,
  control (branch, jump, call/return, trap), system (OS call, virtual memory),
  floating point, decimal, string, graphics. The first four are standard on
  every machine; the others may be primitives, a coprocessor, or synthesized by
  the compiler.
\end{itemize}

\subsubsection{Addressing modes}
\begin{center}
\small
\begin{tabular}{@{}lll@{}}
\toprule
Mode & Example & Operand / effective address \\
\midrule
Immediate & \texttt{ADD R0,\#10} & the constant in the instruction \\
Register & \texttt{ADD R0,R1} & \texttt{R1} \\
Direct (absolute) & \texttt{ADD R0,(100)} & \texttt{M[100]} \\
Register indirect & \texttt{ADD R0,(R1)} & \texttt{M[R1]} \\
Displacement & \texttt{LD R1,100(R2)} & \texttt{M[R2+100]} \\
Indexed & \texttt{ADD R3,(R1+R2)} & \texttt{M[R1+R2]} \\
Memory indirect & \texttt{load r1 = M[M[r2]]} & \texttt{M[M[R2]]} \\
PC-relative & \texttt{branch r1<r3, 1000} & target $=$ PC $+$ offset \\
Scaled & \texttt{load r1 = M[100+r3+r4*d]} & \texttt{M[100+R3+R4$\times$d]} \\
\bottomrule
\end{tabular}
\end{center}
\begin{itemize}[leftmargin=*]
  \item Classified by the \emph{source} of data (immediate, register, memory)
  and by the \emph{address calculation} (direct, indirect, indexed).
  \item SPEC89 on the VAX: \textbf{displacement} (32--55\%) and
  \textbf{immediate} (17--43\%) dominate; register deferred 3--24\%, scaled
  0--16\%, memory indirect 1--6\%. So a new ISA should support these two well.
  \item Most displacements fit in 12--16 bits; small constants are about
  50\% of operands, so a 16-bit immediate field covers most cases.
\end{itemize}

\subsubsection{Operand types and sizes}
Character: ASCII 8-bit (almost always), Unicode 16-bit. Integer (2's
complement): short 16-bit, long 32-bit. Floating point: single 32-bit, double
64-bit. Business: BCD (exact decimal fractions). DSP: fixed point, block
floating point. Graphics: RGBA/XYZW with 8-, 16-bit or single-precision
components.

\subsubsection{Instruction format and registers}
\begin{itemize}[leftmargin=*]
  \item Encoding affects \textbf{program size} and \textbf{decode
  complexity}, and must fit enough registers and addressing modes.
  \item \textbf{Variable length} (VAX, x86): opcode $+$ address specifier/field
  pairs; dense code, hard to decode. \textbf{Fixed length} (ARM, MIPS,
  PowerPC): easy to decode and pipeline, larger code. \textbf{Hybrid}: two
  fixed lengths (MIPS16, Thumb) for density with simple decode.
  \item A field for $n$ choices needs $\lceil\log_2 n\rceil$ bits; $k$-bit
  two's-complement field: $-2^{k-1}\ldots2^{k-1}-1$.
  \item \textbf{Expanding opcode}: leave unused codes of a short opcode as
  escapes that extend into operand fields of shorter instructions.
  \item \textbf{Registers}: fastest storage, flexible for the compiler, but
  cost more per bit; more registers need more bits per instruction and make the
  register file slower.
\end{itemize}

\subsubsection{CISC vs.\ RISC}
\begin{center}
\small
\begin{tabular}{@{}p{3cm}p{5.6cm}p{5.6cm}@{}}
\toprule
& CISC & RISC \\
\midrule
Goal & do more per instruction, fewer instructions & simple instructions, $\CPI\le1$ \\
Motivation & expensive memory (M6800: 16K RAM \$500); backward compatibility & pipelining, compiler synthesizes complex ops \\
Encoding & variable length & fixed length (32 bits), few formats \\
Memory access & ALU operands can be in memory & only load/store \\
Addressing & many modes (14 in MC68000, 25 in MC68020) & immediate, register indirect, displacement, indexed \\
Other & 0- to 3-address instructions (VAX) & delayed loads/branches, prefetch, speculation; more instructions executed \\
Examples & VAX, 680x0, x86 & MIPS, SPARC, PowerPC, Alpha, PA-RISC \\
\bottomrule
\end{tabular}
\end{center}
Case studies: 680x0 has 18 addressing modes and instructions of 2--10 bytes;
x86 has 12 modes and instructions of 1--12 bytes; HP PA-RISC has 7 modes and 12
formats, all 32 bits; Alpha has 4 modes and 7 formats, all 32 bits.

\subsubsection{MIPS}
\textbf{Design principles}: simplicity favors regularity (fixed size, few
formats, opcode always in the first 6 bits); smaller is faster (few
instructions, registers, modes); make the common case fast (register operands,
immediates); good design demands good compromises (three formats).

SPEC2006 instruction mix: integer programs are 35\% data transfer, 34\%
conditional branch, 16\% arithmetic, 12\% logical, 2\% jump; floating-point
programs are 48\% arithmetic, 36\% data transfer, 8\% branch, 4\% logical.

\begin{center}
\small
\begin{tabular}{@{}lcccccc@{}}
\toprule
Format & 31--26 & 25--21 & 20--16 & 15--11 & 10--6 & 5--0 \\
\midrule
R & op (0) & rs & rt & rd & shamt & funct \\
I & op & rs & rt & \multicolumn{3}{c}{immediate (16)} \\
J & op & \multicolumn{5}{c}{address (26)} \\
\bottomrule
\end{tabular}
\end{center}

\begin{center}
\footnotesize
\begin{tabular}{@{}llll@{}}
\toprule
Mode in MIPS & Instructions & Address / operand & Range \\
\midrule
Register & R-format & \texttt{R[rs]}, \texttt{R[rt]} & 32 registers \\
Immediate & \texttt{addi}, \texttt{ori}, \texttt{slti} & 16-bit constant & $-32768\ldots32767$ \\
Base (displacement) & \texttt{lw}, \texttt{sw}, \texttt{lb}, \dots & \texttt{R[rs]} $+$ SignExt(imm) & $\pm32$ KiB \\
PC-relative & \texttt{beq}, \texttt{bne} & PC$+4+$SignExt(imm)$\times4$ & $\pm2^{15}$ instructions \\
Pseudo-direct & \texttt{j}, \texttt{jal} & $\{(\text{PC}{+}4)[31{:}28],\ \text{addr},\ 00\}$ & same 256 MiB region \\
Register indirect & \texttt{jr}, \texttt{jalr} & \texttt{R[rs]} & any address \\
\bottomrule
\end{tabular}
\end{center}

\begin{center}
\small
\begin{tabular}{@{}llll@{}}
\toprule
R-format (op $=0$) & funct & I/J-format & opcode \\
\midrule
\texttt{sll}, \texttt{srl}, \texttt{sra} & 00, 02, 03 & \texttt{j}, \texttt{jal} & 02, 03 \\
\texttt{jr}, \texttt{jalr} & 08, 09 & \texttt{beq}, \texttt{bne}, \texttt{blez}, \texttt{bgtz} & 04, 05, 06, 07 \\
\texttt{add}, \texttt{addu} & 20, 21 & \texttt{addi}, \texttt{addiu}, \texttt{slti}, \texttt{sltiu} & 08, 09, 0a, 0b \\
\texttt{sub}, \texttt{subu} & 22, 23 & \texttt{andi}, \texttt{ori}, \texttt{xori}, \texttt{lui} & 0c, 0d, 0e, 0f \\
\texttt{and}, \texttt{or}, \texttt{xor}, \texttt{nor} & 24, 25, 26, 27 & \texttt{lb}, \texttt{lh}, \texttt{lw}, \texttt{lbu}, \texttt{lhu} & 20, 21, 23, 24, 25 \\
\texttt{slt}, \texttt{sltu} & 2a, 2b & \texttt{sb}, \texttt{sh}, \texttt{sw}; \texttt{ll}, \texttt{sc} & 28, 29, 2b; 30, 38 \\
\bottomrule
\end{tabular}
\end{center}
(All in hex.)

\begin{center}
\small
\begin{tabular}{@{}llll@{}}
\toprule
Name & Number & Use & Preserved across a call? \\
\midrule
\texttt{\$zero} & 0 & constant 0 & n/a \\
\texttt{\$at} & 1 & assembler temporary & no \\
\texttt{\$v0--\$v1} & 2--3 & results & no \\
\texttt{\$a0--\$a3} & 4--7 & arguments & no \\
\texttt{\$t0--\$t7}, \texttt{\$t8--\$t9} & 8--15, 24--25 & temporaries & no \\
\texttt{\$s0--\$s7} & 16--23 & saved & yes \\
\texttt{\$k0--\$k1} & 26--27 & OS kernel & no \\
\texttt{\$gp}, \texttt{\$sp}, \texttt{\$fp}, \texttt{\$ra} & 28--31 & pointers; return address & yes \\
\bottomrule
\end{tabular}
\end{center}
\begin{remark}
\slidefix{} The slide table marks \texttt{\$a0--\$a3} as preserved; the MIPS
reference sheet (H\&P green card) says they are \emph{not}. A callee may
overwrite its arguments, so a caller that still needs them must save them.
\end{remark}

\begin{itemize}[leftmargin=*]
  \item \textbf{R-format}: \texttt{rd $\leftarrow$ rs funct rt}. Shifts use
  \texttt{rt} as the source, \texttt{shamt} as the amount, \texttt{rs}$=0$.
  \texttt{sra} shifts in the sign bit. \texttt{jr rs}: register indirect.
  \item \textbf{I-format}: \texttt{rt} is the destination (or the value stored
  by \texttt{sw}). \textbf{Sign-extended}: \texttt{addi}, \texttt{slti},
  load/store offsets, branch offsets. \textbf{Zero-extended}: \texttt{andi},
  \texttt{ori}, \texttt{xori}.
  \item \texttt{add}, \texttt{addi}, \texttt{sub} trap on signed overflow;
  \texttt{addu}, \texttt{addiu}, \texttt{subu} never trap. \texttt{slt}
  compares signed, \texttt{sltu} unsigned.
  \item \texttt{lb}/\texttt{lh} sign-extend, \texttt{lbu}/\texttt{lhu}
  zero-extend. Words must be 4-byte aligned, halfwords 2-byte aligned.
  \item \texttt{mult}/\texttt{div} write \texttt{HI}/\texttt{LO} (product:
  HI$=$upper 32 bits, LO$=$lower; division: LO$=$quotient, HI$=$remainder);
  read them with \texttt{mfhi}/\texttt{mflo}.
  \item \textbf{Big-endian} (as taught): the most significant byte is at the
  lowest address. Little-endian puts the least significant byte there.
  \item \textbf{Synthesized operations}: no \texttt{MOV}
  (\texttt{add \$t1,\$t0,\$zero}), no \texttt{NOT} (\texttt{nor \$d,\$s,\$zero}),
  no \texttt{SUBI} (\texttt{addi} with a negative constant), no \texttt{SLA}
  (same bits as \texttt{sll}).
  \item \textbf{32-bit constant}: \texttt{lui \$t0,upper; ori \$t0,\$t0,lower}
  (\texttt{ori} because it zero-extends).
  \item \textbf{Atomic read-modify-write}: \texttt{ll} then \texttt{sc} to the
  same address. \texttt{sc} stores and sets \texttt{rt}$=1$ only if the word was
  not written since the \texttt{ll}; otherwise it stores nothing and sets
  \texttt{rt}$=0$, so the code retries.
  \item \textbf{Procedures}: \texttt{jal label} puts the return address in
  \texttt{\$ra} and jumps; \texttt{jr \$ra} returns. Convention: save return
  address, save registers, run, restore registers, restore return address,
  return. The stack grows down: \texttt{addi \$sp,\$sp,-4k} then
  \texttt{sw}, and the reverse on exit. A non-leaf procedure must save
  \texttt{\$ra}.
  \item \textbf{Assembler}: directives \texttt{.text}, \texttt{.data},
  \texttt{.globl}, \texttt{.word}, \texttt{.space}. Pseudo-instructions
  (\texttt{li}, \texttt{la}, \texttt{move}, \texttt{lw} of a label, \texttt{blt},
  \dots) expand to real instructions, using \texttt{\$at} as scratch. SPIM's
  data segment starts at \texttt{0x10010000}.
\end{itemize}

% ----------------------------------------------------------------------------
\subsection{Common traps}
\begin{enumerate}[leftmargin=*]
  \item Branch offsets count \emph{instructions} from PC$+4$, not bytes from
  the branch: offset $=(\text{target}-(\text{PC}+4))/4$.
  \item The \texttt{j} field is the target \emph{divided by 4} (26 bits); the
  top 4 bits come from PC$+4$.
  \item Shifts: the source is \texttt{rt}, not \texttt{rs}; \texttt{rs}$=0$.
  \item I-format: the destination is \texttt{rt}, written \emph{second} in the
  binary but \emph{first} in assembly.
  \item \texttt{ori}/\texttt{andi} zero-extend; \texttt{addi} sign-extends.
  \texttt{lui}+\texttt{addi} breaks when the low half is $\ge$\texttt{0x8000}.
  \item \texttt{lb} vs.\ \texttt{lbu}: a byte $\ge$\texttt{0x80} becomes
  negative with \texttt{lb}.
  \item Code size and memory traffic are different questions: count
  instruction bytes for code size, data accesses for traffic.
  \item Fewer instructions (CISC) does not mean faster: use $I\times\CPI\times C$.
\end{enumerate}

% ----------------------------------------------------------------------------
\subsection{Lecturer's review questions}

\begin{exercise}[Lecturer 1]
In the MIPS instruction set, \texttt{NOT A} can be replaced by:
(a) \texttt{A NAND 0}; (b) \texttt{A XOR 0}; (c) \texttt{A NOR 0}; (d) none of the
above.
\end{exercise}
\begin{solution}
\textbf{(c)}: $\overline{A+0}=\overline{A}$, i.e.\ \texttt{nor \$d,\$s,\$zero}.
\texttt{A NAND 0} is always all ones and \texttt{A XOR 0} $=A$ (and MIPS has no
\texttt{nand}).
\end{solution}

\begin{exercise}[Lecturer 2]
Microprocessor performance can be measured by: (a) MIPS; (b) throughput; (c) MFLOPS;
(d) all of the above.
\end{exercise}
\begin{solution}
\textbf{(d)}. All three are performance \emph{metrics}. Remember the caveat from
Section~\ref{sec:intro}: MIPS and MFLOPS can mislead, and execution time of real
programs is the only reliable measure.
\end{solution}

\begin{exercise}[Lecturer 3]
The MIPS architecture is a: (a) stack; (b) accumulator; (c) register/register;
(d) register/memory architecture.
\end{exercise}
\begin{solution}
\textbf{(c)}: a load-store (register-register) ISA; only \texttt{lw}/\texttt{sw}
(and friends) access memory.
\end{solution}

\begin{exercise}[Lecturer 4]
Find \texttt{\$t1} after \texttt{lui \$t1,0x1234; addi \$t1,\$t1,0x5678}:
(a) \texttt{0x00000000}; (b) \texttt{0x1234}; (c) \texttt{0x5678};
(d) \texttt{0x12345678}.
\end{exercise}
\begin{solution}
\textbf{(d)}. \texttt{lui} gives \texttt{0x12340000}; \texttt{0x5678} has bit 15 $=0$,
so sign extension gives \texttt{0x00005678} and the sum is \texttt{0x12345678}.
(With a low half $\ge$\texttt{0x8000}, \texttt{addi} would subtract
\texttt{0x10000}; see Exercise~\ref{q:isa-sign}(e).)
\end{solution}

\begin{exercise}[Lecturer 5]
The x86 architecture is: (a) big-endian; (b) most-significant byte at the highest
address of a word; (c) most-significant byte at the lowest address of a word;
(d) none of the above.
\end{exercise}
\begin{solution}
\textbf{(b)}: x86 is \emph{little-endian}; the least-significant byte is at the
lowest address, so the most-significant byte is at the highest. (c) describes
big-endian, the MIPS convention in this course.
\end{solution}

% ----------------------------------------------------------------------------
\subsection{Exam-style questions}
Solutions are in Section~\ref{sec:isa-sol}.

\begin{exercise}[Comparing ISA styles]\label{q:isa-styles}
Compile $E=(A-B)\times(C+D)$, where $A$--$E$ are in memory, for a stack,
accumulator, memory-memory (3-operand), register-memory (3-operand, last operand
may be in memory, as in \texttt{ADD R3,R1,B}) and load-store machine. Do not
overwrite $A$--$D$; use memory temporaries \texttt{T}, \texttt{T1},
\texttt{T2} where needed.
Encoding: opcode 8 bits, memory address 16 bits, register field 4 bits (16
registers); stack \texttt{ADD}/\texttt{SUB}/\texttt{MUL} have only an opcode.
\begin{enumerate}[label=(\alph*),leftmargin=*]
  \item Write the code and count the instructions.
  \item Find the code size in bytes.
  \item Count the data memory accesses (reads $+$ writes of $A$--$E$ and
  temporaries). Assuming 4-byte data, find the total memory traffic (code $+$
  data) per execution.
  \item The next statement is $F=A\times B$. Which styles can reuse work
  already done?
\end{enumerate}
\end{exercise}

\begin{exercise}[Addressing modes]\label{q:isa-modes}
Registers: \texttt{R1}$=100$, \texttt{R2}$=200$, \texttt{R3}$=4$. Memory
(word at each address): \texttt{M[100]}$=300$, \texttt{M[104]}$=50$,
\texttt{M[116]}$=9$, \texttt{M[200]}$=100$, \texttt{M[204]}$=7$,
\texttt{M[300]}$=40$. Find the value loaded into \texttt{R4} and name the mode:
\begin{center}
\small
\begin{tabular}{@{}llll@{}}
(a) \texttt{LOAD R4,\#100} & (b) \texttt{LOAD R4,R1} & (c) \texttt{LOAD R4,(100)} & (d) \texttt{LOAD R4,(R2)} \\
(e) \texttt{LOAD R4,4(R2)} & (f) \texttt{LOAD R4,(R1+R3)} & (g) \texttt{LOAD R4,@(R1)} & (h) \texttt{LOAD R4,0(R1)[R3]} \\
\end{tabular}
\end{center}
(\texttt{@} is memory indirect; (h) is scaled: \texttt{M[0+R1+R3$\times$d]} with $d=4$.)
\begin{enumerate}[label=(\alph*),leftmargin=*,start=9]
  \item Which of these modes does MIPS have? Write MIPS code for (c), (f) and
  (g) with \texttt{\$s1}$=$\texttt{R1}, \texttt{\$s3}$=$\texttt{R3},
  \texttt{\$s4}$=$\texttt{R4}.
  \item Using the SPEC89 data, explain why MIPS keeps only a few modes.
\end{enumerate}
\end{exercise}

\begin{exercise}[Instruction format design]\label{q:isa-format}
\leavevmode
\begin{enumerate}[label=(\alph*),leftmargin=*]
  \item A machine has 16-bit instructions and 16 registers. It needs 15
  instructions with three register operands, 14 with two, 31 with one and 16
  with none. Design an expanding-opcode encoding, or show it is impossible.
  \item MIPS is extended to 128 registers, keeping 32-bit instructions and a
  6-bit opcode. How wide can the I-format immediate be now, and what is its
  range? How many bits are left in the R-format for \texttt{shamt} and
  \texttt{funct}? What are the consequences?
  \item Give two reasons for and two against adding more registers.
  \item Why does a fixed-length encoding help pipelining, and how do MIPS16 and
  Thumb recover code density?
\end{enumerate}
\end{exercise}

\begin{exercise}[CISC vs.\ RISC]\label{q:isa-cisc}
A program is compiled for a CISC and a RISC machine:
\begin{center}
\begin{tabular}{@{}lcc@{}}
\toprule
& CISC & RISC \\
\midrule
Static instruction count & 50{,}000 & 80{,}000 \\
Average instruction length & 3.5 bytes & 4 bytes \\
Dynamic instruction count & $1.0\times10^9$ & $1.6\times10^9$ \\
CPI & 4.0 & 1.2 \\
Clock & 2.0 GHz & 2.5 GHz \\
\bottomrule
\end{tabular}
\end{center}
\begin{enumerate}[label=(\alph*),leftmargin=*]
  \item Find the code size on each machine.
  \item Find the CPU time on each machine and the speedup of RISC.
  \item A compressed RISC mode (like Thumb) needs 15\% more instructions;
  70\% of them are 2 bytes and the rest 4 bytes. Find the new code size.
  \item Explain why CISC made sense in the 1970s and why it was kept.
\end{enumerate}
\end{exercise}

\begin{exercise}[Encoding MIPS]\label{q:isa-enc}
Give the format, fields, binary and hex of:
(a) \texttt{slt \$t0,\$s1,\$s2};
(b) \texttt{srl \$t1,\$s3,4};
(c) \texttt{andi \$s2,\$t3,0x00F0};
(d) \texttt{lbu \$t2,3(\$a1)};
(e) \texttt{sw \$ra,4(\$sp)};
(f) \texttt{jal Proc}, with \texttt{Proc} at \texttt{0x00400200}.
\end{exercise}

\begin{exercise}[Decoding MIPS]\label{q:isa-dec}
The words below are at addresses \texttt{0x00400000}, \texttt{0x00400004},
\dots{} in order. Decode each one; for branches and jumps give the target
address.
\begin{center}
\texttt{0x02119822}\quad \texttt{0x2108FFFF}\quad \texttt{0xAFBF0004}\quad
\texttt{0x1520FFFC}\quad \texttt{0x0C100080}
\end{center}
\end{exercise}

\begin{exercise}[Branch and jump ranges]\label{q:isa-range}
\leavevmode
\begin{enumerate}[label=(\alph*),leftmargin=*]
  \item \texttt{beq \$s0,\$s1,Far} is at \texttt{0x00400100} and \texttt{Far}
  is at \texttt{0x00480000}. Can it be encoded? If not, rewrite it and give the
  encoding of the jump.
  \item Can a \texttt{j} at \texttt{0x00400100} reach \texttt{0x10000000}? If
  not, how do you get there?
  \item State the reach of \texttt{beq}, \texttt{j} and \texttt{jr}.
  \item Why are branch and jump offsets stored in words rather than bytes?
\end{enumerate}
\end{exercise}

\begin{exercise}[Bytes, signs and overflow]\label{q:isa-sign}
MIPS is big-endian. The word \texttt{0x8A0B7CFF} is stored at address
\texttt{0x2000}.
\begin{enumerate}[label=(\alph*),leftmargin=*]
  \item Give the byte at each address \texttt{0x2000}--\texttt{0x2003}. Find
  the 32-bit result of \texttt{lb} and \texttt{lbu} from \texttt{0x2000} and
  from \texttt{0x2003}, \texttt{lh} from \texttt{0x2000} and \texttt{lhu} from
  \texttt{0x2002}.
  \item What does \texttt{lb} from \texttt{0x2000} return on a little-endian
  machine? What happens with \texttt{lw} from \texttt{0x2002}?
  \item \texttt{\$t1}$=$\texttt{0x7FFFFFFF}. What do \texttt{addi \$t2,\$t1,1}
  and \texttt{addiu \$t2,\$t1,1} do?
  \item \texttt{\$s0}$=$\texttt{0xFFFFFFFF}, \texttt{\$s1}$=1$. Find
  \texttt{slt \$t0,\$s0,\$s1} and \texttt{sltu \$t0,\$s0,\$s1}.
  \item Load \texttt{0x8000F000} into \texttt{\$t0}. What goes wrong with
  \texttt{lui \$t0,0x8000; addi \$t0,\$t0,0xF000}?
\end{enumerate}
\end{exercise}

\begin{exercise}[Directives and pseudo-instructions]\label{q:isa-asm}
In SPIM the data segment starts at \texttt{0x10010000}.
\begin{lstlisting}
        .data
x:      .word  5
buf:    .space 6
y:      .word  0x11223344
        .text
        .globl main
main:   la    $t0, y
        lw    $t1, x
        li    $t2, 0x12345
        move  $a0, $t1
\end{lstlisting}
\begin{enumerate}[label=(\alph*),leftmargin=*]
  \item Give the addresses of \texttt{x}, \texttt{buf} and \texttt{y}.
  \item Expand each instruction of \texttt{main} into real MIPS instructions.
  \item Why is \texttt{\$at} reserved for the assembler?
\end{enumerate}
\end{exercise}

\begin{exercise}[C to MIPS and instruction counts]\label{q:isa-count}
\leavevmode
\begin{lstlisting}[language=C,keywords={int,char,while,if,return}]
int count(char *s, char c) {      // s in $a0, c in $a1
    int n = 0;
    while (*s != 0) { if (*s == c) n++; s++; }
    return n;                     // in $v0
}
\end{lstlisting}
\begin{enumerate}[label=(\alph*),leftmargin=*]
  \item Write \texttt{count} in MIPS. Does it need the stack?
  \item How many instructions run for \texttt{s = "abca"}, \texttt{c = 'a'}?
  \item ALU instructions take 1 cycle and loads, branches and jumps take 2.
  Find the cycles and the CPI.
  \item Move the loop test to the bottom to remove the \texttt{j}. Repeat (b)
  and (c) and give the speedup.
\end{enumerate}
\end{exercise}

\begin{exercise}[Recursion and the stack]\label{q:isa-fact}
Use the recursive \texttt{fact} from the course:
\begin{lstlisting}
fact: addi $sp, $sp, -8
      sw   $ra, 4($sp)
      sw   $a0, 0($sp)
      slti $t0, $a0, 1
      beq  $t0, $zero, L1
      addi $v0, $zero, 1
      addi $sp, $sp, 8
      jr   $ra
L1:   addi $a0, $a0, -1
      jal  fact
      lw   $a0, 0($sp)
      lw   $ra, 4($sp)
      addi $sp, $sp, 8
      mul  $v0, $a0, $v0
      jr   $ra
\end{lstlisting}
\begin{enumerate}[label=(\alph*),leftmargin=*]
  \item \texttt{main} calls \texttt{fact(3)} with \texttt{\$sp}$=$\texttt{0x7FFFEFFC}.
  What is the lowest value \texttt{\$sp} reaches, and what is stored in each
  frame?
  \item How many instructions of \texttt{fact} run in total?
  \item Give \texttt{\$v0} at each return.
  \item Why must \texttt{\$a0} and \texttt{\$ra} be saved? What happens if
  \texttt{\$ra} is not?
\end{enumerate}
\end{exercise}

\begin{exercise}[Atomic increment with LL/SC]\label{q:isa-llsc}
\leavevmode
\begin{lstlisting}
again: ll   $t0, 0($s0)
       addi $t1, $t0, 1
       sc   $t1, 0($s0)
       beq  $t1, $zero, again
\end{lstlisting}
\begin{enumerate}[label=(\alph*),leftmargin=*]
  \item What does the code do, and what is in \texttt{\$t0} at the end?
  \item Two processors run it at the same time on a counter equal to 5. Trace
  one possible interleaving. What is the final value? What could happen with
  \texttt{lw}/\texttt{sw} instead?
  \item Why should the code between \texttt{ll} and \texttt{sc} be short?
\end{enumerate}
\end{exercise}

\begin{exercise}[Short answers]\label{q:isa-short}
Answer in one or two sentences.
\begin{enumerate}[label=(\alph*),leftmargin=*]
  \item Why does MIPS have no \texttt{SUBI}, no \texttt{NOT} and no
  \texttt{SLA}?
  \item List three uses of \texttt{\$zero}.
  \item Why is the opcode always in bits 31--26 and \texttt{rs}, \texttt{rt}
  always in the same place?
  \item Why is a load-store ISA easier to pipeline than a register-memory one?
  \item Which instruction classes should be fast for SPEC2006 integer and for
  floating-point programs?
  \item Why does MIPS need \texttt{ll}/\texttt{sc} instead of a single atomic
  instruction?
  \item Name an operand type a business machine needs that a DSP does not,
  and the reverse.
  \item Which part of the CPU equation does a RISC ISA increase, and which does
  it reduce?
\end{enumerate}
\end{exercise}

% ----------------------------------------------------------------------------
\subsection{Solutions}
\label{sec:isa-sol}

\begin{noot}{Exercise~\ref{q:isa-styles}}
(a), (b) Instruction sizes: \texttt{PUSH}/\texttt{POP}, accumulator
instructions $8+16=24$ bits; stack ALU $8$; memory-memory $8+3\cdot16=56$;
register load/store $8+4+16=28$; register-memory ALU $8+4+4+16=32$;
register-register ALU $8+3\cdot4=20$.
\begin{center}
\footnotesize
\begin{tabular}{@{}p{1.9cm}p{6.1cm}cp{2.5cm}c@{}}
\toprule
Style & Code & Instr. & Code size & Data acc. \\
\midrule
Stack & \texttt{PUSH A; PUSH B; SUB; PUSH C; PUSH D; ADD; MUL; POP E} & 8 &
$5\cdot24+3\cdot8=144$ b $=18$ B & 5 \\
Accumulator & \texttt{LOAD C; ADD D; STORE T; LOAD A; SUB B; MUL T; STORE E} & 7 &
$7\cdot24=168$ b $=21$ B & 7 \\
Memory-memory & \texttt{SUB T1,A,B; ADD T2,C,D; MUL E,T1,T2} & 3 &
$3\cdot56=168$ b $=21$ B & 9 \\
Register-memory & \texttt{LOAD R1,A; SUB R1,R1,B; LOAD R2,C; ADD R2,R2,D; MUL R3,R1,R2; STORE R3,E} & 6 &
$2(28{+}32){+}20{+}28$ $=168$ b $=21$ B & 5 \\
Load-store & \texttt{LOAD R1,A; LOAD R2,B; LOAD R3,C; LOAD R4,D; SUB R1,R1,R2; ADD R3,R3,R4; MUL R1,R1,R3; STORE R1,E} & 8 &
$5\cdot28+3\cdot20=200$ b $=25$ B & 5 \\
\bottomrule
\end{tabular}
\end{center}
(c) Total traffic $=$ code $+$ $4\times$data accesses: stack $18+20=38$ B,
accumulator $21+28=49$ B, memory-memory $21+36=57$ B, register-memory
$21+20=41$ B, load-store $25+20=45$ B. The stack machine has the densest code;
memory-memory and accumulator machines pay with data traffic (temporaries go to
memory); load-store has the largest code but the simplest instructions.

(d) Only the register styles: \texttt{A} and \texttt{B} are still in
\texttt{R1}/\texttt{R2} on the load-store machine (if the compiler keeps
them) and \texttt{A} in \texttt{R1} on the register-memory one, so no reloads
are needed. Stack, accumulator and memory-memory machines must read \texttt{A}
and \texttt{B} from memory again.
\end{noot}

\begin{noot}{Exercise~\ref{q:isa-modes}}
(a) immediate: 100. (b) register: 100. (c) direct: \texttt{M[100]}$=300$.
(d) register indirect: \texttt{M[200]}$=100$. (e) displacement:
\texttt{M[204]}$=7$. (f) indexed: \texttt{M[104]}$=50$. (g) memory indirect:
\texttt{M[M[100]]}$=$\texttt{M[300]}$=40$. (h) scaled:
\texttt{M[100+4$\times$4]}$=$\texttt{M[116]}$=9$.

(i) MIPS has register, immediate and displacement (plus PC-relative and
pseudo-direct for control). Register indirect is displacement with 0; direct is
displacement from \texttt{\$zero} (addresses below 32 KiB only):
\begin{lstlisting}
lw   $s4, 100($zero)      # (c) direct
add  $t0, $s1, $s3        # (f) indexed: form the address first
lw   $s4, 0($t0)
lw   $t0, 0($s1)          # (g) memory indirect: two loads
lw   $s4, 0($t0)
\end{lstlisting}
(j) Displacement and immediate cover most uses (32--55\% and 17--43\%); memory
indirect ($\le6\%$) and scaled ($\le16\%$) are rare and can be synthesized
with one extra instruction. Fewer modes keep decode simple and the clock short
(smaller is faster, make the common case fast).
\end{noot}

\begin{noot}{Exercise~\ref{q:isa-format}}
\begin{enumerate}[label=(\alph*),leftmargin=*]
  \item Possible, with every code used. Three operands: 4-bit opcode
  \texttt{0000}--\texttt{1110} (15), \texttt{1111} is the escape. Two operands:
  \texttt{1111 cccc rrrr rrrr}, \texttt{cccc} $=$ \texttt{0000}--\texttt{1101}
  (14), escapes \texttt{1110}, \texttt{1111}. One operand:
  \texttt{1111 111x cccc rrrr}, 32 codes, use 31 and keep
  \texttt{1111 1111 1111} as the escape. None: \texttt{1111 1111 1111 cccc},
  16 codes. Check: $15\cdot2^{12}+14\cdot2^8+31\cdot2^4+16=61440+3584+496+16=65536=2^{16}$.
  \item Register fields are 7 bits. I-format: $32-6-7-7=12$ bits, range
  $-2048\ldots2047$. R-format: $32-6-3\cdot7=5$ bits for \texttt{shamt} and
  \texttt{funct} together. Consequences: offsets and constants need
  \texttt{lui}/extra instructions more often, branches reach only $\pm2^{11}$
  instructions, and there is no room for both a 5-bit \texttt{shamt} and a
  \texttt{funct}, so fewer R-type operations (or variable shifts only).
  \item For: fewer loads/stores (spills), more values kept close for the
  compiler. Against: more bits per instruction, a slower (larger) register file,
  more state to save on calls and context switches.
  \item Every instruction has the same length and field positions, so the next
  PC is known (PC$+4$) and decode and register read can start at once, in fixed
  stages. MIPS16/Thumb add a second, 16-bit fixed format for common
  instructions: denser code, still simple decode.
\end{enumerate}
\end{noot}

\begin{noot}{Exercise~\ref{q:isa-cisc}}
\begin{enumerate}[label=(\alph*),leftmargin=*]
  \item CISC $50{,}000\times3.5=175{,}000$ B; RISC $80{,}000\times4=320{,}000$ B
  ($1.83\times$ larger).
  \item CISC $\dfrac{10^9\times4.0}{2\times10^9}=2.0$ s; RISC
  $\dfrac{1.6\times10^9\times1.2}{2.5\times10^9}=0.768$ s. Speedup
  $2.0/0.768=2.6$: 60\% more instructions, but a much lower CPI and faster clock.
  \item $80{,}000\times1.15=92{,}000$ instructions, average
  $0.7\cdot2+0.3\cdot4=2.6$ B, size $239{,}200$ B: 25\% smaller than plain RISC,
  still larger than CISC.
  \item Memory was expensive (16 KiB RAM for \$500, a 40 MB disk \$55{,}000),
  so compact code mattered, and complex instructions narrowed the gap to
  high-level languages. Later CISC designs had to run old binaries, so the
  instruction sets were kept for backward compatibility.
\end{enumerate}
\end{noot}

\begin{noot}{Exercise~\ref{q:isa-enc}}
\begin{center}
\footnotesize
\setlength{\tabcolsep}{3pt}
\begin{tabular}{@{}lp{4.1cm}l@{}}
\toprule
Instruction & Fields & Binary $=$ hex \\
\midrule
(a) R & op 0, rs 17, rt 18, rd 8, sh 0, fn \texttt{2a} &
\texttt{000000 10001 10010 01000 00000 101010} $=$ \texttt{0x0232402A} \\
(b) R & op 0, rs 0, rt 19, rd 9, sh 4, fn 02 &
\texttt{000000 00000 10011 01001 00100 000010} $=$ \texttt{0x00134902} \\
(c) I & op \texttt{0c}, rs 11, rt 18, imm \texttt{0x00F0} &
\texttt{001100 01011 10010 0000000011110000} $=$ \texttt{0x317200F0} \\
(d) I & op 24, rs 5, rt 10, imm 3 &
\texttt{100100 00101 01010 0000000000000011} $=$ \texttt{0x90AA0003} \\
(e) I & op \texttt{2b}, rs 29, rt 31, imm 4 &
\texttt{101011 11101 11111 0000000000000100} $=$ \texttt{0xAFBF0004} \\
(f) J & op 03, addr $=$ \texttt{0x00400200}$/4=$ \texttt{0x0100080} &
\texttt{000011 00000100000000000010000000} $=$ \texttt{0x0C100080} \\
\bottomrule
\end{tabular}
\end{center}
\end{noot}

\begin{noot}{Exercise~\ref{q:isa-dec}}
\begin{itemize}[leftmargin=*]
  \item \texttt{0x02119822} $=$ \texttt{000000 10000 10001 10011 00000 100010}:
  R, funct \texttt{22} $\Rightarrow$ \texttt{sub \$s3,\$s0,\$s1}.
  \item \texttt{0x2108FFFF} $=$ \texttt{001000 01000 01000 1111111111111111}:
  \texttt{addi \$t0,\$t0,-1}.
  \item \texttt{0xAFBF0004} $=$ \texttt{101011 11101 11111 0000000000000100}:
  \texttt{sw \$ra,4(\$sp)}.
  \item \texttt{0x1520FFFC} $=$ \texttt{000101 01001 00000 1111111111111100}:
  \texttt{bne \$t1,\$zero,-4}. At \texttt{0x0040000C}: target
  $=$ \texttt{0x00400010}$-16=$ \texttt{0x00400000}.
  \item \texttt{0x0C100080} $=$ \texttt{000011} $+$ \texttt{0x0100080}:
  \texttt{jal}, target $=\{0000,\ \texttt{0x0100080}\times4\}=$ \texttt{0x00400200}.
\end{itemize}
\end{noot}

\begin{noot}{Exercise~\ref{q:isa-range}}
\begin{enumerate}[label=(\alph*),leftmargin=*]
  \item Offset $=(\texttt{0x00480000}-\texttt{0x00400104})/4=131{,}007>32{,}767$:
  it does not fit. Invert the condition and jump:
\begin{lstlisting}
0x00400100        bne  $s0, $s1, Skip   # offset 1
0x00400104        j    Far              # 0x08120000
0x00400108  Skip: ...
\end{lstlisting}
  \texttt{j} field $=$ \texttt{0x00480000}$/4=$ \texttt{0x0120000}, encoding
  \texttt{0x08120000}; the top 4 bits of PC$+4$ and of \texttt{Far} are both 0.
  \item No: the top 4 bits of PC$+4$ are \texttt{0000} but the target's are
  \texttt{0001}. Use \texttt{lui \$at,0x1000; jr \$at}.
  \item \texttt{beq}: $-2^{15}\ldots2^{15}-1$ instructions from PC$+4$, about
  $\pm128$ KiB. \texttt{j}: anywhere in the current 256 MiB region.
  \texttt{jr}: any 32-bit address.
  \item Instructions are word-aligned, so the low 2 bits of every target are 0;
  storing words gives $4\times$ the reach for free.
\end{enumerate}
\end{noot}

\begin{noot}{Exercise~\ref{q:isa-sign}}
\begin{enumerate}[label=(\alph*),leftmargin=*]
  \item \texttt{0x2000}: \texttt{8A}, \texttt{0x2001}: \texttt{0B},
  \texttt{0x2002}: \texttt{7C}, \texttt{0x2003}: \texttt{FF}.
  \texttt{lb 0x2000} $=$ \texttt{0xFFFFFF8A} ($-118$);
  \texttt{lbu 0x2000} $=$ \texttt{0x0000008A} (138);
  \texttt{lb 0x2003} $=$ \texttt{0xFFFFFFFF} ($-1$);
  \texttt{lbu 0x2003} $=$ \texttt{0x000000FF} (255);
  \texttt{lh 0x2000} $=$ \texttt{0xFFFF8A0B};
  \texttt{lhu 0x2002} $=$ \texttt{0x00007CFF}.
  \item Little-endian puts \texttt{FF} at \texttt{0x2000}, so \texttt{lb}
  returns \texttt{0xFFFFFFFF}. \texttt{lw} from \texttt{0x2002} is not
  word-aligned: address error exception.
  \item \texttt{addi}: $2^{31}-1+1$ overflows, so an exception is raised and
  \texttt{\$t2} is not written. \texttt{addiu}: \texttt{\$t2}$=$\texttt{0x80000000}, no exception.
  \item \texttt{slt}: $-1<1$, so 1. \texttt{sltu}: $4{,}294{,}967{,}295<1$ is false, so 0.
  \item \texttt{lui \$t0,0x8000; ori \$t0,\$t0,0xF000}. With \texttt{addi},
  \texttt{0xF000} is sign-extended to \texttt{0xFFFFF000}; adding it to
  \texttt{0x80000000} ($-2^{31}$) gives signed overflow (an exception), and even
  \texttt{addiu} would give \texttt{0x7FFFF000}.
\end{enumerate}
\end{noot}

\begin{noot}{Exercise~\ref{q:isa-asm}}
\begin{enumerate}[label=(\alph*),leftmargin=*]
  \item \texttt{x} $=$ \texttt{0x10010000} (4 B); \texttt{buf} $=$
  \texttt{0x10010004}--\texttt{0x10010009}; \texttt{.word} aligns to 4, so
  \texttt{y} $=$ \texttt{0x1001000C} (2 bytes of padding).
  \item One valid expansion:
\begin{lstlisting}
lui  $at, 0x1001          # la $t0, y
ori  $t0, $at, 0x000C
lui  $at, 0x1001          # lw $t1, x
lw   $t1, 0($at)
lui  $at, 0x0001          # li $t2, 0x12345
ori  $t2, $at, 0x2345
addu $a0, $zero, $t1      # move $a0, $t1
\end{lstlisting}
  \item Pseudo-instructions need a scratch register to build addresses and
  constants; if programs used \texttt{\$at}, the expansion would destroy their
  values.
\end{enumerate}
\end{noot}

\begin{noot}{Exercise~\ref{q:isa-count}}
(a) A leaf procedure that only uses \texttt{\$t} registers needs no stack:
\begin{lstlisting}
count: add  $v0, $zero, $zero   # n = 0
loop:  lbu  $t0, 0($a0)         # *s
       beq  $t0, $zero, done
       bne  $t0, $a1, next
       addi $v0, $v0, 1         # n++
next:  addi $a0, $a0, 1         # s++ (chars are bytes)
       j    loop
done:  jr   $ra
\end{lstlisting}
(b) Each character: \texttt{lbu}, \texttt{beq}, \texttt{bne}, \texttt{addi},
\texttt{j} $=5$, plus 1 on a match. Total $1+4\cdot5+2+2+1=26$ (2 matches,
final \texttt{lbu}+\texttt{beq}, \texttt{jr}).

(c) ALU 7 (\texttt{add}, 2 \texttt{addi} $n$, 4 \texttt{addi} $s$), loads 5,
branches 9, jumps 5. Cycles $=7+2(5+9+5)=45$, $\CPI=45/26=1.73$.

(d)
\begin{lstlisting}
count: add  $v0, $zero, $zero
       lbu  $t0, 0($a0)
       beq  $t0, $zero, done
loop:  bne  $t0, $a1, next
       addi $v0, $v0, 1
next:  addi $a0, $a0, 1
       lbu  $t0, 0($a0)
       bne  $t0, $zero, loop
done:  jr   $ra
\end{lstlisting}
Instructions $=3+4\cdot4+2+1=22$: ALU 7, loads 5, branches 9, jumps 1.
Cycles $=7+2(5+9+1)=37$, $\CPI=37/22=1.68$. Speedup $=45/37=1.22$.
\end{noot}

\begin{noot}{Exercise~\ref{q:isa-fact}}
\begin{enumerate}[label=(\alph*),leftmargin=*]
  \item Four calls (\texttt{n}$=3,2,1,0$), 8 bytes each: lowest
  \texttt{\$sp} $=$ \texttt{0x7FFFEFFC}$-32=$ \texttt{0x7FFFEFDC}. Each frame
  holds \texttt{\$ra} at \texttt{4(\$sp)} and \texttt{\$a0} at
  \texttt{0(\$sp)}: frame 1 has \texttt{main}'s return address and 3; frames
  2--4 have the address after \texttt{jal fact} and 2, 1, 0.
  \item A call with \texttt{n}$\ge1$ runs 5 (entry and test) $+2$
  (\texttt{addi}, \texttt{jal}) $+5$ (after return) $=12$; the base case runs
  $5+3=8$. Total $3\cdot12+8=44$.
  \item \texttt{fact(0)}$=1$, \texttt{fact(1)}$=1\cdot1=1$,
  \texttt{fact(2)}$=2$, \texttt{fact(3)}$=6$.
  \item \texttt{\$a0} is changed for the recursive call but \texttt{n} is
  needed for the multiply afterwards; \texttt{jal} overwrites \texttt{\$ra}.
  Without saving \texttt{\$ra}, every \texttt{jr \$ra} returns to the
  instruction after \texttt{jal fact}, so the function never returns to
  \texttt{main} (it loops).
\end{enumerate}
\end{noot}

\begin{noot}{Exercise~\ref{q:isa-llsc}}
\begin{enumerate}[label=(\alph*),leftmargin=*]
  \item An atomic fetch-and-increment of \texttt{M[\$s0]}: \texttt{ll} reads
  the value, \texttt{sc} writes value$+1$ only if nobody wrote the word
  meanwhile, and otherwise the loop retries. \texttt{\$t0} holds the old value.
  \item P1 \texttt{ll} (5), P2 \texttt{ll} (5), P1 \texttt{sc} 6 succeeds, P2
  \texttt{sc} fails (the word changed) and returns 0, P2 retries: \texttt{ll}
  (6), \texttt{sc} 7. Final value 7. With \texttt{lw}/\texttt{sw} both could
  read 5 and write 6: one increment is lost.
  \item The longer the window, the more likely another write (or an interrupt)
  makes \texttt{sc} fail, so the loop may retry many times.
\end{enumerate}
\end{noot}

\begin{noot}{Exercise~\ref{q:isa-short}}
\begin{enumerate}[label=(\alph*),leftmargin=*]
  \item \texttt{SUBI} is \texttt{addi} with a negative constant; \texttt{NOT}
  is \texttt{nor \$d,\$s,\$zero}; a left arithmetic shift produces the same bits
  as \texttt{sll}. Leaving them out keeps the instruction set small.
  \item \texttt{move} (\texttt{add \$d,\$s,\$zero}), clear a register,
  compare with zero (\texttt{beq \$t0,\$zero,L}), negate
  (\texttt{sub \$d,\$zero,\$s}), direct addressing (\texttt{lw \$t0,100(\$zero)}).
  \item Regularity: the hardware can read \texttt{rs}/\texttt{rt} from the
  register file while it decodes the opcode, so decode is simple and fast.
  \item Only loads and stores touch memory, so every instruction uses each
  stage the same way (one memory access at most, in a fixed stage) and has a
  predictable length and timing.
  \item Integer: data transfers (35\%) and conditional branches (34\%).
  FP: arithmetic (48\%) and data transfers (36\%).
  \item A read-modify-write in one instruction needs both a read and a write
  in one memory operation, which breaks the load-store pipeline; the pair keeps
  each instruction simple and still guarantees atomicity.
  \item Business: BCD for exact decimal money values. DSP: fixed point or block
  floating point.
  \item It raises the instruction count $I$ and lowers the CPI (and allows a
  shorter clock cycle $C$).
\end{enumerate}
\end{noot}
